\documentclass[12pt]{article}
\usepackage[utf8]{inputenc}
\usepackage[T1]{fontenc}
\usepackage[margin=1in]{geometry}
\usepackage{setspace}
\usepackage{xcolor}
\usepackage[colorlinks=true, urlcolor=blue!60!black, linkcolor=blue!60!black]{hyperref}
\setlength{\parindent}{0em}
\setlength{\parskip}{0.45em}
\renewcommand{\baselinestretch}{1.0}
\begin{document}
\begin{center}
{\Large \textit{United States v. Nixon}, 418 U.S. 683 (1974)\par}
\vspace{0.3em}
{\normalsize Case Brief --- One-Page Edition\par}
\end{center}

\textbf{Memory Jogger.} The Watergate tapes: executive privilege is real but qualified (\href{https://supreme.justia.com/cases/federal/us/418/683/}{418 U.S. 683}).

\textbf{Facts.} Special Prosecutor Leon Jaworski subpoenaed tapes of 64 Oval Office conversations for the trial of seven Nixon aides; the grand jury had named Nixon an unindicted coconspirator. Nixon claimed absolute privilege.

\textbf{Procedural History.} Judge Sirica ordered \textit{in camera} production; the Court granted certiorari before judgment. Nixon resigned two weeks after the decision.

\textbf{Votes.} Unanimous, 8--0; Burger wrote; Rehnquist did not participate.

\textbf{Holding.} Presidential communications are presumptively privileged, but absent military or diplomatic secrets a generalized claim of confidentiality yields to the demonstrated, specific need for evidence in a pending criminal trial.

\textbf{Reasoning.} The appeal was proper and the dispute justiciable because the Executive was bound by its own Special Prosecutor regulation. Courts, not the President, decide the scope of privilege. The privilege is rooted in Article II, but an absolute privilege would impair the courts and the guarantees of due process and compulsory process.

\textbf{Examples.} \emph{Same-Side:} a specific subpoena for tapes of meetings about a bribe. \emph{Opposite-Side:} a subpoena for classified war plans. \emph{Fence-Sitter:} a congressional subpoena for White House memoranda.

\textbf{Critique.} Originalists doubt a constitutional privilege exists; others note the Court armed later presidents against oversight.

\textbf{Key Quotation.} ``A President and those who assist him must be free to explore alternatives in the process of shaping policies and making decisions, and to do so in a way many would be unwilling to express except privately. These are the considerations justifying a presumptive privilege for Presidential communications. The privilege is fundamental to the operation of Government, and inextricably rooted in the separation of powers under the Constitution.'' --- Chief Justice Burger, Opinion of the Court, 418 U.S.\ at 708.

\end{document}
